% Préparation d'une solution fille par dilution (corrigé de l'exercice 5, chapitre 2)
\begin{tikzpicture}[x=1cm,y=1cm,
  etape/.style={font=\otLight\fontsize{7.5}{9.2}\selectfont,text=otInk,align=center,text width=33mm,anchor=north}]
% 1. Prélèvement
\pic at (0,0) {becher={largeur=1.5,hauteur=1.8,niveau=1.0,liquide=cyan!35}};
\pic at (0.05,0.25) {pipette={hauteur=4.1,niveau=1,liquide=cyan!35}};
\node[pcLegendeG] at (-0.45,0.55) {S$_0$};
\node[etape] at (0,-0.35) {\otRond{1}\enspace Prélever le volume $V_0$ de solution mère S$_0$ avec une pipette jaugée munie d'une propipette.};
% 2. Transfert
\pic at (3.7,0) {fiole={hauteur=3.1,niveau=0.18,liquide=cyan!35}};
\pic at (3.7,2.55) {pipette={hauteur=2.9,niveau=0.35,liquide=cyan!35}};
\node[etape] at (3.7,-0.35) {\otRond{2}\enspace Verser dans une fiole jaugée de volume $V_1$.};
% 3. Complément et homogénéisation
\pic at (7.3,0) {fiole={hauteur=3.1,niveau=1,liquide=cyan!15}};
\pic at (6.05,1.25) {pissette};
\draw[pcRepere] (7.54,2.35) -- (8.05,2.75) node[pcLegende,anchor=west]{trait de jauge};
\node[pcLegendeG] at (7.3,0.55) {S$_1$};
\node[etape] at (7.15,-0.35) {\otRond{3}\enspace Compléter avec de l'eau distillée jusqu'au trait de jauge, boucher, homogénéiser.};
\end{tikzpicture}
\par\smallskip\legende{Ici, $V_0=\qty{10,0}{mL}$ de sirop et $V_1=\qty{100,0}{mL}$ : le facteur de dilution vaut 10. Schéma sans échelle.}
